\section{数学推理}

\answer{AG2-1}
扣除已知对角边在顶点 $1,3$ 的贡献，剩余四个校验都是零。每个顶点的两条环边于是必须相等，沿环传递后，四条环边全零或全一，恰有两解。四个公平独立位共有 $16$ 种等可能赋值，所以概率为 $2/16=1/8$。再观察到 $12=0$ 后只剩全零解；未知边构成一条路径，也符合 A19 的森林唯一性条件。

\answer{AG2-2}
列与两种对角线的斜率分别为 $a,a+1,1-a$，在模 $5$ 下均须非零。因此排除 $a=0,1,4$，留下 $a=2,3$；对每个合法斜率，$b$ 任取五个余数，共 $10$ 种排列。$9$ 被 $3$ 整除，A14 的平方和必要条件已不成立，所以不存在任何满足要求的排列，不能通过改用非仿射方式绕过障碍。

\answer{AG2-3}
跳跃右端为 $3/8+1/4=5/8$，所以条件分布为 $\operatorname{Unif}(3/8,5/8)$。其中大于 $1/2$ 的部分占一半，因此联合概率为 $(1/4)(1/2)=1/8$。
原随机变量只在可数集合 $\{r_j\}$ 取值，故 $h(X)$ 几乎处处落在可数集合 $\{h(r_j)\}$。连续均匀分布对任何可数集合的概率都为零，与此集合概率为一矛盾。即使原子位置在实数轴上稠密，也不改变这个结论。

\answer{AG2-4}
被删成员原来持有三种物品，它们各从三人共享变为两人共享；其余四种仍由三人持有。所以两人种数为 $3$、三人种数为 $4$，没有物品消失。任何剩余成员持有的物品没被删除，仍有三种；任何剩余人员对原来共同的唯一物品也仍在。
另令两人、三人种数为 $x,y$。持有次数与共享人员对数分别给出 $2x+3y=6\cdot3=18$、$x+3y=\binom62=15$。相减得 $x=3$，继而 $y=4$，与构造一致。

\answer{AG3-1}
藏点均匀时每次搜查命中概率为 $1/2$，搜查对均匀时每个藏点被覆盖的概率也为 $1/2$，故值为 $1/2$。
若藏点概率为 $p_i$，四种搜查的命中概率平均为 $1/2$。要使其最大值不超过 $1/2$，必须全部等于 $1/2$，即 $p_1=p_3=a$、$p_2=p_4=1/2-a$，其中 $0\le a\le1/2$。
对搜查方，四个藏点的覆盖率平均同样为 $1/2$。要使最小覆盖率至少为 $1/2$，必须各点覆盖率都为 $1/2$。以 $q_{ij}$ 表示搜查该边的概率，全部解为
\[
 q_{12}=q_{34}=b,\qquad q_{23}=q_{41}=1/2-b,
 \qquad0\le b\le1/2.
\]
这些条件既必要又充分，因此包括区间端点在内给出全部最优策略。

\answer{AG3-2}
$1,2,3$ 颗分别为败、胜、败：只有一颗必须取完而输，两颗可留一颗，三颗可取一颗留给对方胜态或取完而输。$4,5,6,7$ 颗可分别取 $3,4,3,4$ 颗，留下 $1$ 或 $3$。所以前七个正局面只有 $1,3$ 为败。
一般必败余数为 $1,3\pmod7$。由这两个余数取 $1,3,4$ 不会到另一个必败余数；若动作恰好取完，则本来就是输的动作。其余余数 $0,2,4,5,6$ 可分别取 $4,1,3,4,3$ 颗，留下正的必败局面；最小正代表处也合法。强归纳给出全部分类。
$20\equiv6\pmod7$，新规则应取 $3$，留下 $17\equiv3$。另两种动作留下的 $19,16$ 都是胜态，所以这是唯一必胜首步。A09 的正常规则下应取 $4$，留下 $16\equiv2$，同样是唯一必胜首步。

\answer{AG3-3}
两边均有人时，不偏离条件变为
\[
 k\le\frac{9-k}{2}+\delta,\qquad
 \frac{8-k}{2}+\delta\le k+1,
\]
也就是 $6+2\delta\le3k\le9+2\delta$。要使 $k=3$ 成为均衡，需 $0\le\delta\le3/2$。在此范围内空队配置都不是均衡；当 $0<\delta<3/2$ 时唯一整数解为 $k=3$。
在 $\delta=0$ 时，$k=2$ 也是均衡；在 $\delta=3/2$ 时，$k=4$ 也是均衡。范围外 $k=3$ 本身不满足偏离条件。这里不会出现没有纯均衡的空集情形：当 $0\le\delta<15/2$ 时，取 $k=\lceil2+2\delta/3\rceil$ 即满足两边均有人时的均衡条件；当 $\delta\ge15/2$ 时，$k=8$ 为均衡。因此所求恰为 $0<\delta<3/2$。这个人数配置的总服务成本仍为 $43/2$；费用改变激励，而不改变给定人数下的实际服务成本。

\answer{AG3-4}
顾客在中点 $(a+b)/2$ 处分流，分段积分得平均距离
\[
 C(a,b)=\frac{a^2}{2}+\frac{(b-a)^2}{4}
       +\frac{(1-b)^2}{2}.
\]
梯度为 $((3a-b)/2,(3b-a-2)/2)$，唯一驻点是 $a=1/4,b=3/4$。Hessian 的特征值为 $1,2$，所以目标严格凸；驻点在可行区域内，故是唯一有序全局最优，代入得 $C=1/8$。
竞争时，不同位置的左方可向右靠近对方以严格增加需求，所以均衡必须重合。重合点必须是均匀需求的中位数 $1/2$，否则向较多需求的一侧偏移可得超过一半的份额；重合中点时任何偏离都不会获益。因此唯一纯均衡为 $(1/2,1/2)$，平均距离为 $\int_0^1|t-1/2|\,\dd t=1/4$，是合作最优的两倍。

\par\Needspace{8\baselineskip}
\subsection*{小节练习}
\answer{OA17}
设总共有 $N$ 种，第 $j$ 种出现在 $r_j$ 个人的收藏中，其中 $1\le r_j\le4$。计数所有人持有的种类数，得到
\[
 \sum_{j=1}^N r_j=40.
\]
再把六对收藏的交集大小相加。第 $j$ 种被计入 $\binom{r_j}{2}$ 次，因此
\[
 \sum_{j=1}^N\binom{r_j}{2}\le36,
 \qquad \sum_{j=1}^N r_j^2\le112.
\]
Cauchy--Schwarz 不等式给出
\[
 40^2\le N\sum_{j=1}^N r_j^2\le112N,
 \qquad N\ge\frac{100}{7}.
\]
$N$ 是整数，故 $N\ge15$。

把明信片类型编号为 $1$ 至 $15$，取以下四个收藏：
\[
 \begin{aligned}
 A&=\{1,2,3,4,5,6,7,8,9,10\},\\
 B&=\{1,2,3,4,5,6,11,12,13,14\},\\
 C&=\{1,2,7,8,9,10,11,12,13,14\},\\
 D&=\{3,4,5,7,8,9,11,12,13,15\}.
 \end{aligned}
\]
每个集合都有十个元素，六个两两交集均恰有六个元素，并集包含全部十五种。因此最小值为 $\boxed{15}$。

\answer{OA22}
若 $x>y$，则 $3x-2y=x+2(x-y)\ge4$。若 $x=y$，整数条件排除 $x=1$，故 $3x-2y=x\ge2$。以下只需考虑 $y=x+d$、$d\ge1$。

令
\[
 P=(x+1)(x+2)\cdots(x+d).
\]
原分式等于 $x!P/(P+1)$。因为 $P$ 与 $P+1$ 互素，整数条件等价于 $P+1\mid x!$。

反设结论不成立，即 $x-2d\le1$，所以 $x\le2d+1$。任取 $P+1$ 的素因子 $p$。由整除关系，$p\le x$。若 $p\le d$，连续的 $d$ 个整数 $x+1,\ldots,x+d$ 中至少有一个是 $p$ 的倍数，导致 $p\mid P$，矛盾。因此 $p>d$。若还满足 $p\le y/2$，则
\[
 x<2d+2\le2p\le y,
\]
于是 $2p$ 是 $P$ 的一个因子，又矛盾。因此所有素因子均满足
\[
 \frac y2<p\le x.
\]
若 $d\ge x$，这个区间为空，与 $P+1>1$ 矛盾。否则其中的整数个数为
\[
 x-\left\lfloor\frac{x+d}{2}\right\rfloor
 =\left\lceil\frac{x-d}{2}\right\rceil\le d.
\]
最后一个不等式来自 $x\le2d+1$ 和 $d\ge1$。又因 $p>y/2>x/2$，素数 $p$ 在 $x!$ 中只出现一次。由 $P+1\mid x!$，$P+1$ 只能是至多 $d$ 个不同的、不超过 $x$ 的素数之积。于是
\[
 P+1\le x^d<(x+1)\cdots(x+d)=P,
\]
矛盾。故 $3x-2y\ge2$。取 $(x,y)=(4,5)$ 时原分式为 $20$，且 $3x-2y=2$，说明常数不能提高。

\answer{OA28}
总权重为 $\sum_{n=0}^\infty2^{-n}=2$，所以级数在每个 $x$ 处均收敛，且任意尾部的贡献都可统一控制。

先设 $x$ 为无理数。给定 $\varepsilon>0$，选 $N$ 使 $\sum_{n>N}2^{-n}<\varepsilon$。有限多个点 $r_0,\ldots,r_N$ 均不等于 $x$，故存在 $\delta>0$，使 $|y-x|<\delta$ 时这些点与 $y$ 的大小关系和与 $x$ 的大小关系相同。前 $N+1$ 项不变，于是
\[
 |f(y)-f(x)|\le\sum_{n>N}2^{-n}<\varepsilon.
\]
所以每个无理点都连续。

再设 $x=r_k$。分离第 $k$ 项，写成
\[
 f(y)=2^{-k}\mathbf1_{\{y>r_k\}}+g(y).
\]
对 $g$ 使用同样的有限前缀加尾部估计：前缀中已没有等于 $r_k$ 的点，故 $g$ 在 $r_k$ 连续。严格不等式使 $f(r_k)=g(r_k)$，从而
\[
 \boxed{\lim_{y\uparrow r_k}f(y)=f(r_k),\qquad
 \lim_{y\downarrow r_k}f(y)=f(r_k)+2^{-k}.}
\]
因此每个有理点都有大小为 $2^{-k}$ 的向上跳跃，左连续而不连续；全部连续点恰为无理数。跳点虽稠密，但尾部总权重很小，故不会破坏无理点的连续性。

\answer{OA09}
对正整数 $r$，用 $\ell(r)$ 表示整除 $r$ 的最大的 $2$ 的幂，即二进制中最低非零位的值。先研究一般局面 $(r,m)$：还剩 $r$ 颗，本次最多取 $m$ 颗。下面证明
\[
 (r,m)\text{ 为必败局面}\quad\Longleftrightarrow\quad m<\ell(r).
\]
证明对剩余石子数归纳。若 $m\ge\ell(r)$，取走 $\ell(r)$ 颗。若恰好取完则获胜；否则 $r-\ell(r)$ 的最低非零位严格高于 $\ell(r)$，而对方的取数上限为 $\ell(r)$，故对方进入必败局面。

若 $m<\ell(r)$，任取合法数量 $a$，都有 $1\le a<\ell(r)$。写 $a=2^s q$，其中 $q$ 为奇数。因 $r$ 是更大的 $2$ 的幂的倍数，$r-a$ 恰含 $s$ 个因子 $2$，所以
\[
 \ell(r-a)=\ell(a)\le a.
\]
对方的上限为 $a$，故进入必胜局面。所有合法走法都把必胜局面交给对手，完成归纳。

初始局面为 $(n,n-1)$，它必败当且仅当 $\ell(n)=n$，也就是 $n$ 为$2$ 的幂。因此
\[
 \boxed{\text{先手必胜}\iff n\notin\{2^k:k\ge1\}.}
\]
必胜时第一步取 $\ell(n)$ 颗；以后每当轮到自己，就取当前剩余数量的最低非零位。上述判据保证这一步始终合法，并最终取完石子。

\answer{OA13}
先手第一步选择 $6$。此后所有合法数都不能含素因子 $2$ 或 $3$，剩下可能出现的素因子只有
\[
 5,7,11,13,17,19,23,29.
\]
任意两个不同的这些素数的乘积至少为 $5\cdot7=35>30$。所以以后每个合法数恰好含其中一种素因子；例如 $25$ 只使用素因子 $5$。选过一个数后，它使用的素因子就不能再出现。

因而每步恰好消耗上述八个素因子中的一个。只要还有未使用的素因子 $p$，数 $p$ 本身仍是合法选择，游戏就不会提前结束。因此在首步 $6$ 之后必定还有八步，与双方之后的具体选择无关。

全局共走九步，先手走最后一步，对手随后无数可选。因此先手选择 $6$ 后即可保证获胜。证明的关键是既说明每步消耗一个素因子，也说明尚有素因子时一定还有合法数。

